\documentclass[11pt]{article}

\usepackage[a4paper,margin=0.88in]{geometry}
\usepackage{amsmath,amssymb,amsthm}
\usepackage{graphicx}
\usepackage[hidelinks]{hyperref}
\usepackage{microtype}

\newtheorem{theorem}{Theorem}
\newtheorem{corollary}{Corollary}
\newcommand{\ellp}[1]{\ell^{#1}}

\title{Exact Mixed-Norm Bernstein Numbers\\in the Aligned Parameter Regions}
\author{Automatic math-discovery run \texttt{fa\_banach\_001}}
\date{August 11, 2026}

\begin{document}
\maketitle

\section*{Source Question}

The source is Chuah--Lang--Yao, \emph{A Classification Theorem on
Non-compact Embeddings between Besov Spaces}, arXiv:2410.10731
\cite{CLY}.  For
\[
  0<p_0<p_1\leq\infty,\qquad 0<q_0<q_1\leq\infty,
\]
Proposition~3.5 proves that the formal embedding
\[
 I:\ellp{q_0}(\ellp{p_0})\longrightarrow
       \ellp{q_1}(\ellp{p_1})
\]
satisfies
\begin{equation}\label{eq:source-upper}
 b_n(I)\leq
 n^{-\min(1/p_0,1/q_0)
       (1-\max(p_0/p_1,q_0/q_1))}.
\end{equation}
Section~6 asks for sharper upper bounds and matching lower bounds, and
singles out the decay of these mixed-sequence Bernstein numbers as the
critical question for noncompact Besov embeddings on $\mathbb R^d$.
More precisely, the source writes:
\begin{quote}
``We believe that the upper bounds in Proposition 3.5 could be sharpened,
and that it should be possible to derive corresponding lower bounds.''
It then identifies the case $\Omega=\mathbb R^d$ with computing the decay
of $b_n(\ell^{q_0}(\ell^{p_0})\to\ell^{q_1}(\ell^{p_1}))$.
\end{quote}

\begin{center}
\includegraphics[width=0.95\linewidth]{figures/open_problem_crop.png}
\end{center}

Here $1/\infty=0$ and $r/\infty=0$.  We use the source convention
\[
 b_n(T)=\sup_{\dim E=n}\ \inf_{0\ne x\in E}
             \frac{\lVert Tx\rVert}{\lVert x\rVert}.
\]

\section*{Proof Intuition}

There are two canonical $n$-dimensional tests.  Matrices supported in one
row see only the inner $p$-embedding; matrices supported in one column see
only the outer $q$-embedding.  They give a lower bound governed by the
slower of the two scalar $\ell^r$ decay rates.  The source upper exponent is
the product of two minima.  When the smaller initial reciprocal exponent
and the smaller relative loss occur in the same coordinate direction,
that product is exactly the slower scalar rate, so the upper and lower
bounds coincide without any loss.  In crossed regions the two minima come
from different directions; this argument no longer closes the gap, and
those regions remain open in this packet.

\section*{Result}

\begin{theorem}\label{thm:main}
Suppose, in addition, that either
\begin{align}
 &p_0\ge q_0 \quad\text{and}\quad
   \frac{p_0}{p_1}\ge\frac{q_0}{q_1}, \label{eq:cone1}\\
 \intertext{or}
 &p_0\le q_0 \quad\text{and}\quad
   \frac{p_0}{p_1}\le\frac{q_0}{q_1}. \label{eq:cone2}
\end{align}
Then, for every $n\ge1$,
\begin{equation}\label{eq:exact}
 \boxed{\quad
 b_n(I)=n^{-\min\{,1/p_0-1/p_1,\,
                         1/q_0-1/q_1,\}}.
 \quad}
\end{equation}
Thus the source upper bound is already sharp, with constant one, throughout
the two aligned parameter regions \eqref{eq:cone1}--\eqref{eq:cone2}.
\end{theorem}

The alignment says that the smaller of $1/p_0,1/q_0$ and the smaller of
the two relative exponent losses come from the same coordinate direction.
It includes, for example, all cases $p_0=q_0$ and all cases
$p_0/p_1=q_0/q_1$.

\section*{Proof}

Put
\[
 a=\frac1{p_0},\qquad c=\frac1{q_0},\qquad
 u=1-\frac{p_0}{p_1},\qquad
 v=1-\frac{q_0}{q_1}.
\]
The two one-direction coordinate subspaces give the complete lower-bound
argument.  Fix one outer index and let $E_p$ consist of matrices supported
on the first $n$ inner coordinates of that row.  For $x\in E_p$ the outer
norms disappear, and the standard finite-dimensional norm comparison gives
\[
 \frac{\lVert x\rVert_{\ellp{q_1}(\ellp{p_1})}}
      {\lVert x\rVert_{\ellp{q_0}(\ellp{p_0})}}
 =\frac{\lVert x\rVert_{p_1}}{\lVert x\rVert_{p_0}}
 \ge n^{1/p_1-1/p_0}=n^{-au}.
\]
Equality is attained by a constant-modulus vector, so the minimum modulus
on $E_p$ is exactly $n^{-au}$.  Consequently $b_n(I)\ge n^{-au}$.

Next fix one inner coordinate and use the first $n$ outer indices.  The
same argument with $q_0,q_1$ gives $b_n(I)\ge n^{-cv}$.  Therefore
\begin{equation}\label{eq:lower}
 b_n(I)\ge \max\{n^{-au},n^{-cv}\}
          =n^{-\min(au,cv)}.
\end{equation}

The exponent in the source upper bound \eqref{eq:source-upper} is
\begin{equation}\label{eq:paperexp}
 \min(a,c)\bigl(1-\max(1-u,1-v)\bigr)
 =\min(a,c)\min(u,v).
\end{equation}
Under \eqref{eq:cone1}, one has $a\le c$ and $u\le v$, hence
\[
 \min(a,c)\min(u,v)=au=\min(au,cv).
\]
Under \eqref{eq:cone2}, one has $c\le a$ and $v\le u$, hence both sides
equal $cv$.  Thus \eqref{eq:source-upper} and \eqref{eq:lower} coincide,
which proves \eqref{eq:exact}. \qed

\section*{Besov Consequence}

\begin{corollary}
Consider a critical noncompact embedding on $\mathbb R^d$ covered by the
source wavelet reduction, with $p_0<p_1$, $q_0<q_1$, and
\[
 s_0-\frac d{p_0}=s_1-\frac d{p_1}.
\]
If \eqref{eq:cone1} or \eqref{eq:cone2} holds, then its Bernstein numbers
have the exact power-law order
\[
 b_n\asymp
 n^{-\min\{1/p_0-1/p_1,\,1/q_0-1/q_1\}}.
\]
The implicit constants come only from the wavelet isomorphisms defining
the equivalent Besov norms.
\end{corollary}

\begin{proof}
The source identifies both Besov spaces with weighted mixed-sequence
spaces.  At the displayed critical relation the same diagonal weight
occurs on both sides and is removed by a common isometry.  The embedding is
therefore conjugate, up to the fixed wavelet isomorphisms, to $I$ in
Theorem~\ref{thm:main}.  Bernstein numbers are stable under composition
with isomorphisms, giving the assertion.
\end{proof}

\section*{Why the Crossed Regions Remain Open}

If, say, $a<c$ but $u>v$, then
\[
 \min(a,c)\min(u,v)=av
 <\min(au,cv).
\]
Thus the source upper estimate and the coordinate lower bound have a real
gap.  Six focused upgrade attempts were made: retaining the row-count
profile in the repeated-maximum lemma; a row/column projection-rank
dichotomy; transposition via mixed-norm Minkowski inequalities;
factorization through intermediate mixed spaces; flattening through a
single $\ell^r$; and rectangular/random-sign subspaces.  None closes the
gap.

The obstruction is dimension allocation.  A repeated-maximum vector may
hide near-maximal coordinates in an unrecorded row, while for an entire
$n$-dimensional subspace that row has limited projection rank.  A lossless
rank/kernel recursion could plausibly prove the full coordinate exponent,
but the elementary iteration loses dimension and decay simultaneously.
Accordingly, this packet claims only the aligned theorem.

\section*{Verification and Novelty Status}

The proof is an exact exponent comparison plus two coordinate restrictions;
a rational-grid script checks the algebra independently.  The source TeX,
run indexes, exact-formula searches, and mixed-norm/Bernstein-number searches
were checked.  No later full solution or prior statement of this equality
was found in that bounded search.  The status is \emph{candidate new
substantial partial result}, not an exhaustive priority claim.

\section*{Human-Review Recommendation}

Review the finite $\ell^p$ minimum modulus, the alignment equivalences, and
the transfer through the source wavelet isomorphisms.  No claim is made in
either crossed region.

\begin{thebibliography}{9}
\bibitem{CLY}
C. Y. Chuah, J. Lang, and L. Yao,
\emph{A Classification Theorem on Non-compact Embeddings between Besov
Spaces}, arXiv:2410.10731.
\end{thebibliography}

\end{document}
